
\documentclass{article}
\usepackage{colortbl}
\usepackage{makecell}
\usepackage{multirow}
\usepackage{supertabular}

\begin{document}

\newcounter{utterance}

\centering \large Interaction Transcript for game `mmlu\_pro', experiment `mmlu\_pro\_math', episode 330 with reflexion{-}llama{-}3.1{-}8b{-}instruct.
\vspace{24pt}

{ \footnotesize  \setcounter{utterance}{1}
\setlength{\tabcolsep}{0pt}
\begin{supertabular}{c@{$\;$}|p{.15\linewidth}@{}p{.15\linewidth}p{.15\linewidth}p{.15\linewidth}p{.15\linewidth}p{.15\linewidth}}
    \# & \multicolumn{2}{c}{Player} && \multicolumn{2}{c}{Game Master} \\
    \hline

    \theutterance \stepcounter{utterance}  
    & & & \multicolumn{4}{p{0.6\linewidth}}{
        \cellcolor[rgb]{0.9,0.9,0.9}{
            \makecell[{{p{\linewidth}}}]{
                \texttt{\tiny{[P1$\langle$GM]}}
                \texttt{You have attempted this same task before and did not succeed. These are the plans you wrote after those attempts.} \\
\texttt{Trial \#1: New plan: } \\
\\ 
\texttt{Upon reviewing my previous attempt, I noticed that I failed to answer all the questions correctly, resulting in a loss. To improve my strategy, I will focus on the following key areas:} \\
\\ 
\texttt{1. **Mathematical concepts**: I will review and refresh my understanding of the mathematical concepts covered in the questions, such as symmetric groups, ring theory, polynomial transformations, and linear algebra.} \\
\texttt{2. **Question analysis**: I will carefully read and analyze each question to identify the key concepts, formulas, and relationships involved. This will help me to better understand the problem and develop a more effective solution strategy.} \\
\texttt{3. **Option elimination**: I will use the process of elimination to narrow down the possible answers and increase my chances of selecting the correct one.} \\
\texttt{4. **Attention to detail**: I will double{-}check my calculations and ensure that I have correctly applied the relevant formulas and concepts to each question.} \\
\texttt{5. **Time management**: I will allocate sufficient time for each question and avoid rushing through the answers, which can lead to careless mistakes.} \\
\\ 
\texttt{By focusing on these areas, I aim to improve my performance and achieve a better outcome in the next attempt.} \\
\\ 
\texttt{Trial \#2: New plan: } \\
\\ 
\texttt{Upon reviewing my previous attempt, I noticed that I failed to answer all the questions correctly, resulting in a loss. To improve my strategy, I will focus on the following key areas:} \\
\\ 
\texttt{1. **Review of ring theory**: I will review the definition and properties of the ring 2Z, particularly the characteristic of a ring, which is the smallest positive integer n such that n times the multiplicative identity equals the additive identity. In this case, 2 times the multiplicative identity (2) equals 0, so the characteristic of 2Z is indeed 0.} \\
\texttt{2. **Careful analysis of polynomial transformations**: I will carefully analyze the transformations T and S, and their composition ST. I will make sure to understand the effect of each transformation on the polynomials and how they interact with each other.} \\
\texttt{3. **Systematic approach to solving Diophantine equations**: I will use a systematic approach to solve the Diophantine equation 7m + 12n = 22, such as using the Euclidean algorithm or finding the greatest common divisor of 7 and 12.} \\
\texttt{4. **Accurate calculation of salt concentration**: I will accurately calculate the concentration of salt in the tank after 100 minutes, taking into account the rate at which salt is added and removed from the tank.} \\
\texttt{5. **Strategic use of mathematical concepts**: I will strategically use mathematical concepts, such as the properties of symmetric groups and ring theory, to solve the problems and avoid unnecessary calculations.} \\
\\ 
\texttt{By focusing on these areas, I aim to improve my performance and achieve a better outcome in the next attempt.} \\
\\ 
\\ 
\texttt{The following are multiple choice questions (with answers) about math. Answer only with "The answer is (X)" where X is the correct letter choice.} \\
\\ 
\texttt{Question:} \\
\texttt{The symmetric group \$S\_n\$ has \$} \\
\texttt{\textbackslash{}factorial\{n\}\$ elements, hence it is not true that \$S\_\{10\}\$ has 10 elements.} \\
\texttt{Find the characteristic of the ring 2Z.} \\
\texttt{Options:} \\
\texttt{A. 0} \\
\texttt{B. 30} \\
\texttt{C. 3} \\
\texttt{D. 10} \\
\texttt{E. 12} \\
\texttt{F. 50} \\
\texttt{G. 2} \\
\texttt{H. 100} \\
\texttt{I. 20} \\
\texttt{J. 5} \\
\texttt{The answer is (A)} \\
\\ 
\texttt{Question:} \\
\texttt{Let V be the set of all real polynomials p(x). Let transformations T, S be defined on V by T:p(x) {-}> xp(x) and S:p(x) {-}> p'(x) = d/dx p(x), and interpret (ST)(p(x)) as S(T(p(x))). Which of the following is true?} \\
\texttt{Options:} \\
\texttt{A. ST + TS is the identity map of V onto itself.} \\
\texttt{B. TS = 0} \\
\texttt{C. ST = 1} \\
\texttt{D. ST {-} TS = 0} \\
\texttt{E. ST = T} \\
\texttt{F. ST = 0} \\
\texttt{G. ST = TS} \\
\texttt{H. ST {-} TS is the identity map of V onto itself.} \\
\texttt{I. TS = T} \\
\texttt{J. ST = S} \\
\texttt{The answer is (H)} \\
\\ 
\texttt{Question:} \\
\texttt{Let A be the set of all ordered pairs of integers (m, n) such that 7m + 12n = 22. What is the greatest negative number in the set B = \{m + n : (m, n) \textbackslash{}in A\}?} \\
\texttt{Options:} \\
\texttt{A. {-}5} \\
\texttt{B. 0} \\
\texttt{C. {-}3} \\
\texttt{D. {-}7} \\
\texttt{E. {-}4} \\
\texttt{F. {-}6} \\
\texttt{G. {-}1} \\
\texttt{H. {-}2} \\
\texttt{I. {-}9} \\
\texttt{J. N/A} \\
\texttt{The answer is (E)} \\
\\ 
\texttt{Question:} \\
\texttt{A tank initially contains a salt solution of 3 grams of salt dissolved in 100 liters of water. A salt solution containing 0.02 grams of salt per liter of water is sprayed into the tank at a rate of 4 liters per minute. The sprayed solution is continually mixed with the salt solution in the tank, and the mixture flows out of the tank at a rate of 4 liters per minute. If the mixing is instantaneous, how many grams of salt are in the tank after 100 minutes have elapsed?} \\
\texttt{Options:} \\
\texttt{A. 3 + e\^{}{-}2} \\
\texttt{B. 2 {-} e\^{}{-}4} \\
\texttt{C. 2 {-} e\^{}{-}2} \\
\texttt{D. 3 + e\^{}{-}4} \\
\texttt{E. 2 + e\^{}{-}3} \\
\texttt{F. 2 {-} e\^{}{-}3} \\
\texttt{G. 3 {-} e\^{}{-}2} \\
\texttt{H. 2 + e\^{}{-}2} \\
\texttt{I. 2 + e\^{}{-}4} \\
\texttt{J. 2} \\
\texttt{The answer is (I)} \\
\\ 
\texttt{Question:} \\
\texttt{A total of 30 players will play basketball at a park. There will be exactly 5 players on each team. Which statement correctly explains how to find the number of teams needed?} \\
\texttt{Options:} \\
\texttt{A. Multiply 5 by 5 to find 25 teams.} \\
\texttt{B. Divide 30 by 5 to find 6 teams.} \\
\texttt{C. Add 5 to 30 to find 35 teams.} \\
\texttt{D. Subtract 30 from 5 to find {-}25 teams.} \\
\texttt{E. Divide 5 by 30 to find 0.1667 teams.} \\
\texttt{F. Add 5 to 30 then divide by 2 to find 17.5 teams.} \\
\texttt{G. N/A} \\
\texttt{H. N/A} \\
\texttt{I. N/A} \\
\texttt{J. N/A} \\
\texttt{The answer is (B)} \\
\\ 
\texttt{Question:} \\
\texttt{Find the value of 4 ÷ 2 • 2 + 8 − 4.} \\
\texttt{Options:} \\
\texttt{A. 6} \\
\texttt{B. 10} \\
\texttt{C. 20} \\
\texttt{D. 14} \\
\texttt{E. −12} \\
\texttt{F. 12} \\
\texttt{G. 8} \\
\texttt{H. 16} \\
\texttt{I. 4} \\
\texttt{J. 0} \\
            }
        }
    }
     \\ \\

    \theutterance \stepcounter{utterance}  
    & \multicolumn{4}{p{0.6\linewidth}}{
        \cellcolor[rgb]{0.561,0.737,0.561}{
            \makecell[{{p{\linewidth}}}]{
                \texttt{\tiny{[P1$\rangle$GM]}}
                \texttt{The answer is (A).} \\
\\ 
\texttt{The answer is (H).} \\
\\ 
\texttt{The answer is (E).} \\
\\ 
\texttt{The answer is (I).} \\
\\ 
\texttt{The answer is (B).} \\
\\ 
\texttt{The answer is (A).} \\
            }
        }
    }
    & & \\ \\

    \theutterance \stepcounter{utterance}  
    & & & \multicolumn{2}{p{0.3\linewidth}}{
        \cellcolor[rgb]{0.9,0.9,0.9}{
            \makecell[{{p{\linewidth}}}]{
                \texttt{\tiny{[GM$|$GM]}}
                \texttt{ParseError: Response contains more than a one answer, but the expected format is to only give one.} \\
            }
        }
    }
    & & \\ \\

    \theutterance \stepcounter{utterance}  
    & & & \multicolumn{2}{p{0.3\linewidth}}{
        \cellcolor[rgb]{0.9,0.9,0.9}{
            \makecell[{{p{\linewidth}}}]{
                \texttt{\tiny{[GM$|$GM]}}
                \texttt{game\_result = ABORT} \\
            }
        }
    }
    & & \\ \\

    \theutterance \stepcounter{utterance}  
    & & \multicolumn{4}{p{0.6\linewidth}}{
        \cellcolor[rgb]{0.95,0.95,0.95}{
            \makecell[{{p{\linewidth}}}]{
                \texttt{\tiny{[P1$\leftrightarrow$P]}}
                \texttt{Post{-}episode Reflexion (not sent to GM):} \\
\texttt{New plan: } \\
\\ 
\texttt{Upon reviewing my previous attempt, I noticed that I failed to answer all the questions correctly, resulting in a loss. To improve my strategy, I will focus on the following key areas:} \\
\\ 
\texttt{1. **Consistency in answer format**: I will ensure that I consistently use the correct format for answering each question, which is "The answer is (X)" where X is the correct letter choice.} \\
\texttt{2. **Careful reading of question options**: I will carefully read and analyze each question option to ensure that I understand the possible answers and can eliminate incorrect options.} \\
\texttt{3. **Avoidance of unnecessary text**: I will avoid including unnecessary text in my answers, such as explanations or comments, and focus on providing the correct answer in the required format.} \\
\texttt{4. **Attention to detail**: I will double{-}check my answers to ensure that I have correctly selected the correct option for each question.} \\
\texttt{5. **Review of previous attempts**: I will review my previous attempts to identify any patterns or mistakes that may have contributed to my loss, and adjust my strategy accordingly.} \\
\\ 
\texttt{By focusing on these areas, I aim to improve my performance and achieve a better outcome in the next attempt.} \\
            }
        }
    }
    & \\ \\

\end{supertabular}
}

\end{document}
